% Pursuing Stacks (cotes 134-2 to 134-8): archive cote 134-3, site page 4 (PDF page 5), typescript page 4.
% typescript. Draft compared: ps2.tex:148-209.
% AI-transcribed by Claude Opus 5.5 (claude-opus-5-5, adaptive thinking, effort high), Message Batches API, batch ps1 (msgbatch_01Rv5hCaNTtUt4JHx1yqoHPL), 2026-11-12; not yet checked by eye.
\scanpage{134-3/4}
\margin{4}
\mspage{4}

xes allow description of homology and cohomology groups.  As a matter of
fact,  the description is formally analoguous,  and nearly identi\unsure{d|c}al, to
the description of the homology groups of a chain complex  -  and it
would seem therefore that stacks (more specifically,  Gr-stacks) are
in a sense the \uline{closest possible} \inserted[typed]{\uline{non commutative}} \uline{generalization of}  \uline{chain complexes},
the homology groups of the chain complex becoming the homotopy groups
of the ``non commutative chain complex'' or stack.

It is well understood,  since Dold-Puppe\hand{${}^{(3)}$,} that chain complexes
form a category equivalent to the category of abelian group objects in
the category of semi-simplicial sets,  or equivalently, \inserted{to} the category of
semi-simplicial abelian groups.  By this equivalence,  the homology
groups of the chain complex are identified with the \uline{homotopy} groups of
the corresponding homotopy type.  As for the \xout{cxxxxx} homology and cohomolo%
\hand{-}gy groups of this homotopy type,  their description in terms of the
chain complex we started with is kind of delicate (I forgot all about it
I am affraid !).  A fortiori, when a homotopy type is described in terms
of a stack,  i.e. a ``non commutative chain complex'',  there is no immedi-%
ate way for describing its homology or cohomology groups  in terms of
this structure\hand{(*)}. What probably should be done,  is to define \inserted{first} a \uline{nerve}
\margin{3.5. This is an error\unsure{;} it appeared soon after that the cohomology can be expressed much in the same way as for \unsure{simplicial} or \unsure{cubical} complexes, using the ``nerve'' and ``forgetting'' \unsure{the structure} of the stack (i.e. \unsure{part} of the ``primitive'' structure.}
functor  from stacks to semi-simplicial sets  (generalizing the familiar
nerve functor defined o\over{f}{n} the category of categories),  and define homolo-%
gy invariants of a stack via those of the associated semi-simplicial set
(directly suited for calculating these).

\margin{16.}
These reflections on the proper place of the notion of a stack with-%
in standard homotopy algebra  are largely a posteriori  -  the clues the\hand{y}
give are surely not so strong as to \inserted[typed]{give an} \xout{\unsure{be}} imperative feeling for the need
of developping this new approach to the homotopy category.  Rather,  the
imperative feeling comes from the \xout{needs of} intuitions tied up with non
commutative  cohomological algebra over topological spaces,  and more ge\unsure{n}%
\hand{-}nerally over topoi, in the spirit of Giraud's thesis,  where a suitable
formalism for non commutative H$^{i}$'s for  i=0,1 or 2 is developped.  He
develops in extenso the notion of stacks,  we should rather say now
1-stacks,  over a topos,  constantly eluding (and for very understandable
reasons !) the notion of a \xout{two} 2-stack,  appearing closely on the heels
of the 1-stacks.  Keeping in mind that 0-stacks are just ordinary
sheaves of sets\hand{,}  on the space or the topos considered,  the hierarchy o\hand{f}
increasingly higher and more sophisticated notions of 0-stacks, 1-stacks,
2-stacks etc,  which will have to be developped over an arbitrary topos,
just parallel\hand{s} the hierarchy of corresponding notions over the one-point
